\section{Standard alternating cores with every larger tensor count}
\label{sec:standard-all-tensors}

\paragraph{Status, contribution and exact scope.}
This is a separate new derivation.  It leaves HD-AST-07--09 unchanged.
The pair-selection argument and the all-size source summation below are
new here.  They use the jointly derived HD-AST-07 standard two-factor
estimate and the weighted-cycle argument of D-ATT-08; the complete
weighted-cycle proof is included with that attribution.  Author checks
and independent whole-proof reviews are recorded separately.  The complete root reading and /root/g34\_eightfolds whole-proof
audit have passed. Shared weighted-cycle authorship is explicit;
the separate exact implementations are not labelled independent
discoveries of that common lemma.

Let $n\ge15$, put $a=n-1$, and write $U=\operatorname{Std}_n$
for the deleted rational permutation module.  Fix $t\ge3$ and
$V=U^{\otimes t}$, of degree $e=a^t$.  A faithful finite actual
inducing image $J$ is assumed to contain the ordinary normal core
$A_n^t$ on $V$ and to act transitively on its $t$ factors.  Its
character field $F$ contains the original imaginary quadratic field
$K$, which satisfies
\[
                    K\ne\mathbf Q(\sqrt{-3}).                 \tag{AST3.1}
\]
There is no assumption that $F=K$.  For a diagonalizable operator
$T$ of degree $v$ put
$\nu(T)=v-\max_z\dim\ker(T-zI)$.  Thus every scalar multiple
has actual moving rank at least $\nu(T)$.  For $x\in S_n$ put
$k_x=\operatorname{ind}(x)=n-c_x$, where $c_x$ counts all cycles,
including fixed points.

The fully reviewed normalizer argument D-AST-01--02 and
HD-AST-04 gives a unique decomposition
\[
             g=\lambda(g)T(g),\qquad
 T(g)\in\mathcal R_{n,t},\qquad \lambda(g)\in\mu_M,              \tag{AST3.2}
\]
where $\mathcal R_{n,t}$ is the actual rational tensor-wreath image
of $S_n\wr S_t$.  Both coordinates are homomorphisms.  We recall
why the uniqueness matters here.  The automorphisms of each $A_n$
factor are implemented by $S_n$ in this range; after implementing
them and the factor permutation, Schur's lemma leaves only a scalar.
A scalar in the rational wreath image centralizes every individual
represented $A_n$ factor, so it has trivial factor permutation.
Each remaining $S_n$ element then centralizes $A_n$ and is the
identity.  The same argument proves faithfulness of the rational
wreath image.  Thus forgetting $\lambda$ gives an actual
permutation representation on $nt$ letters in $t$ blocks; it is
not an inertia-by-inertia choice of representatives.

\subsection*{HD-AST-10. Pair selection at a factor-fixing branch}

\begin{lemma}[HD-AST-10]
Let a genuine inertia $g$ fix all tensor factors.  Write
$T(g)=U(x_1)\otimes\cdots\otimes U(x_t)$ and
$I_g=\sum_{j=1}^t k_{x_j}$.  Then, for its actual scalar phase,
\[
 \operatorname{rank}(1-g)\ge
 a^{t-2}\left(\frac{n-3}{t}I_g+\frac n2\right).                \tag{AST3.3}
\]
\end{lemma}

\begin{proof}
Suppose first that $I_g>0$.  The average of $k_{x_i}+k_{x_j}$
over the $\binom t2$ unordered pairs is $2I_g/t$.  Choose a pair
whose sum is at least that average; its sum is positive.
HD-AST-07, which applies to every target eigenvalue, gives
\[
 \nu\bigl(U(x_i)\otimes U(x_j)\bigr)
 \ge\frac{(n-3)(k_{x_i}+k_{x_j})+n}{2}
 \ge\frac{n-3}{t}I_g+\frac n2.                               \tag{AST3.4}
\]
If a diagonalizable degree-$v$ operator has largest eigenspace
dimension $m$, its tensor product with a degree-$w$ operator has
each eigenspace of dimension at most $mw$: sum the first
multiplicities against the second multiplicities, whose total is
$w$.  Applying this convolution bound to the other $t-2$ factors
multiplies the support in (AST3.4) by $a^{t-2}$.  The actual
scalar can only select one of these target eigenspaces, proving
(AST3.3).

If $I_g=0$, all $x_j$ are the identity.  Since the inertia is
genuine, $g$ is a nonidentity scalar and has moving rank $e$.
The right side of (AST3.3) is then $a^{t-2}n/2\le e$, since
$(n-1)^2\ge n/2$.  This proves the remaining case.  In particular
the full constant $n/2$ is retained even if only one factor acts
nontrivially; it is not divided by the number of tensor factors.
\end{proof}

\subsection*{HD-AST-11. The exact weighted-cycle error}

\begin{lemma}[HD-AST-11, the D-ATT-08 weighted-cycle argument]
Let $X$ be an invertible diagonalizable degree-$a$ operator, and
let $T$ be a weighted tensor $t$-cycle with cycle product $X$.
Define the exact integer
\[
                 R_t(a)=\sum_{j=1}^{t-1}a^{\gcd(t,j)}.          \tag{AST3.5}
\]
For every nonzero scalar $\lambda$,
\[
 \operatorname{rank}(1-\lambda T)\ge
 \left(1-\frac1t\right)a^t+\frac{a^{t-1}}t\nu(X)
                          -\frac{R_t(a)}t.                    \tag{AST3.6}
\]
For $a\ge4$ and $t\ge3$,
\[
                          R_t(a)\le2a^{t-2}.                  \tag{AST3.7}
\]
\end{lemma}

\begin{proof}
Successive changes of basis in the tensor factors conjugate $T$
to a cyclic rotation with $X$ on one factor and identities on the
others.  Explicitly, for
$(v_1,\ldots,v_t)\mapsto(X_1v_t,X_2v_1,\ldots,X_tv_{t-1})$,
use change-of-basis factors
$I,X_2,X_3X_2,\ldots,X_t\cdots X_2$.  The resulting cycle
product is $X_1X_t\cdots X_2$.  Choosing another starting factor
conjugates that product and does not change its eigenspaces.
In this normal form $T^t=X^{\otimes t}$.

Choose an eigenbasis with eigenvalues $\alpha_1,\ldots,\alpha_a$,
repeated with their multiplicities.  For a proposed eigenvalue $z$
of $T$, let $N_z$ count the ordered words $(i_1,\ldots,i_t)$
with $\prod_j\alpha_{i_j}=z^t$.  Fixing the first $t-1$
indices leaves at most $a-\nu(X)$ possibilities for the last,
so
\[
                         N_z\le a^{t-1}(a-\nu(X)).             \tag{AST3.8}
\]
A word orbit under cyclic rotation has a length $d$ dividing $t$.
Its span is a weighted cyclic matrix of degree $d$, with nonzero
weights.  Its characteristic polynomial is $Z^d-c$ with $c\ne0$,
so each of its eigenvalues occurs exactly once.  If it contributes
to the $z$-eigenspace, all its $d$ words are counted by $N_z$.
Its contribution to $t\dim\ker(T-zI)-N_z$ is then $t-d$;
an orbit which does not contribute has nonpositive contribution.
Full-length orbits therefore require no error term.

Let $C_d(a)$ be the number of cyclic word orbits of exact length
$d$ for $d\mid t$.  The total possible short-orbit error is
\[
 \sum_{\substack{d\mid t\\d<t}}(t-d)C_d(a)=R_t(a).            \tag{AST3.9}
\]
To prove the equality, a word orbit of length $d$ has $d$ words,
each fixed by $t/d-1$ nonidentity rotations.  Its contribution to
the count of pairs consisting of a word and a nonidentity rotation
fixing it is $d(t/d-1)=t-d$.  Rotation by $j$ positions fixes
exactly $a^{\gcd(t,j)}$ words.  Counting the same pairs in these
two ways proves (AST3.9), including the factor $t-d$.
Consequently $t\dim\ker(T-zI)\le N_z+R_t(a)$, and (AST3.8)
gives (AST3.6).  Repeated eigenvalues do not merge the distinct
basis-word orbits, so the proof includes all such degeneracies.

For completeness the universal bound (AST3.7) has a short direct
proof from (AST3.5).  For $t=3,4,5$ the errors are respectively
$2a$, $a^2+2a$ and $4a$, each at most $2a^{t-2}$ for $a\ge4$.
For $t\ge6$, each $\gcd(t,j)\le\lfloor t/2\rfloor$, hence
\[
 R_t(a)\le(t-1)a^{\lfloor t/2\rfloor}.
 \qquad
 t-1\le2a^{\lceil t/2\rceil-2}.                              \tag{AST3.10}
\]
The second inequality holds at $t=6$ and $t=7$ for $a\ge4$.
Increasing $t$ by two multiplies its right side by $a$, while
$t+1\le a(t-1)$.  Induction in both parities proves it for every
$t\ge6$, and substitution proves (AST3.7).  This proof is not an
extrapolation from a finite list of tensor counts.
\end{proof}

\subsection*{HD-AST-12. Exclusion for all larger tensor counts}

\begin{theorem}[HD-AST-12]
In the setting (AST3.1)--(AST3.2), no actual curve-cover source has
six-dimensional multiplicity variation with period-image rank nine.
Thus every transitive ordinary standard core $A_n^t$ with
$n\ge15$ and $t\ge3$ is excluded for each non-Eisenstein original
imaginary quadratic field.  Together with HD-AST-09 this excludes
all such tensor counts $t\ge2$ over those fields.  HD-AST-09
also retains its stronger all-field conclusion at $t=2$.
\end{theorem}

\begin{proof}
The actual intermediate-cover factorization of D-AST-03 and
HD-AST-04--09 applies without identifying $F$ with $K$.
If $h$ is the intermediate base genus and $D$ the sum of the
actual moving ranks at its genuine marks, the multiplicity Euler
identity is $6=2e(h-1)+D$.  Since $e>6$, genus at least two
is impossible.  Genus one has $D=6$ and at most six genuine
marks; finite cover data over the pointed base have dimension
at most six, so the actual period image cannot have rank nine.
The proposed rank-nine case must consequently have
\[
 h=0,\qquad r\ge12,\qquad D=2e+6,                            \tag{AST3.11}
\]
for a genuine generating product-one inertia tuple in $J$.
The inequality $r\ge12$ follows from the actual Hurwitz family
over the $r$ pointed sphere, of dimension $r-3$.  No sufficiency
of this dimension condition is being asserted.

Rationality of $T(g)$ gives $F\subseteq\mathbf Q(\zeta_M)$.
If all scalar values lay in $\mu_6$, it would force
$K=\mathbf Q(\sqrt{-3})$, contrary to (AST3.1).  The images of
the tuple therefore generate a nontrivial quotient
$\lambda(J)/(\lambda(J)\cap\mu_6)$ and have product one.
At least two entries have $\lambda\notin\mu_6$.

We recall the exact support estimate HD-AST-05 at these marks.
Every power of a rational tensor-wreath operator has trace at
least $-e/a$.  Indeed its trace is a product of standard
cycle-product traces, each between $-1$ and $a$; a negative
product has a factor $-1$, and there are at most $t$ factors.
For a rational matrix, a primitive $d$th root eigenspace has
equal multiplicity at all its cyclotomic conjugates.  If
$d\notin\{1,2,3,4,6\}$, its multiplicity is at most $e/4$.
For $d=4$, the equal $i,-i$ multiplicities $m$ satisfy
$-e/a\le\operatorname{tr}(T^2)\le e-4m$.
Thus each of the two scalar marks above has actual moving rank
at least
\[
                           \frac{3-1/a}{4}e.                  \tag{AST3.12}
\]
This is applied to the selected eigenspace of the rational
operator before its scalar twist; it is not a claim that a
nonreal-eigenvalue bound survives arbitrary scalar convolution.

Every genuine entry has actual moving rank at least $e/a$.
For a factor-fixing nonscalar entry, an active standard factor
has projective support at least one, and convolution gives this
bound.  A genuine scalar has full moving rank.  If an entry moves
the factors, select one tensor cycle of length $l\ge2$ and put
it in the weighted normal form used above.  Of the $a^l$ basis
words, $a$ are constant and all other rotation orbits have length
at least two.  Each orbit contributes at most one vector to a
specified eigenspace.  Its largest eigenspace fraction is therefore
at most $(1+a^{1-l})/2\le(1+1/a)/2$.  Tensoring with the other
cycles preserves this bound.  Every factor-moving entry thus has
actual moving rank at least
\[
                          s e,\qquad s=\frac{a-1}{2a}\ge\frac1a.
                                                                    \tag{AST3.13}
\]

If three or more entries moved factors, choose the two scalar
marks of (AST3.12) and a third, distinct factor-moving entry.
The remaining at least nine genuine entries give
\[
 D\ge\left(2\frac{3-1/a}{4}+\frac{a-1}{2a}
                         +\frac9a\right)e
       =2e+\frac{8e}{a}>2e+6,                               \tag{AST3.14}
\]
contrary to (AST3.11).  There are at most two factor-moving
entries.  Transitivity requires at least one, and product one
forbids exactly one.  Hence there are precisely two.  Their
factor permutations are inverse, and their cyclic subgroup is
transitive, so both are full $t$-cycles.  The same estimate would
force the two scalar marks to be these two moving entries, but
no use of their scalar orders is made in the rest of the proof.

Let $h_1,h_2\in S_n$ be the products around those two full
cycles, and put $H=k_{h_1}+k_{h_2}$.  Every full-letter
eigenspace of $h_i$ has dimension at most $c_{h_i}$; deleting
one trivial line cannot increase it.  Therefore
\[
                 \nu(U(h_i))\ge a-c_{h_i}=k_{h_i}-1.         \tag{AST3.15}
\]
This intentionally weaker inequality remains valid when $h_i=1$.
Applying HD-AST-11 and (AST3.7), the sum of the two moving
ranks is at least
\[
 2\left(1-\frac1t\right)e+\frac{a^{t-1}}tH
                  -\frac{2a^{t-1}}t-\frac{4a^{t-2}}t.        \tag{AST3.16}
\]

Let $N=\sum I_g$ over the factor-fixing entries.  The actual
permutation representation in (AST3.2) is transitive on $nt$
letters: the core is transitive within each block, and the
factor action is transitive on the blocks.  The inverse image of
a letter stabilizer in the actual connected cover defines a
connected degree-$nt$ quotient cover of the same sphere.
A factor-fixing entry has index $I_g$ there.  A full tensor
cycle with product $h_i$ has index $n(t-1)+k_{h_i}$, because
each length-$l$ cycle of $h_i$ gives a length-$tl$ cycle on
the letters.  Riemann--Hurwitz on this actual quotient implies
\[
               N+H+2n(t-1)\ge2nt-2,
               \qquad N+H\ge2n-2=2a.                       \tag{AST3.17}
\]
Scalar marks have zero permutation index but remain genuine
coefficient-cover marks in the moving-rank sum.

There are at least ten factor-fixing entries.  Sum (AST3.3)
and (AST3.16), then use $a-(n-3)=2>0$ and (AST3.17):
\[
\begin{split}
 D&\ge2\left(1-\frac1t\right)e
       +\frac{a^{t-2}(n-3)}tN+\frac{a^{t-1}}tH
       +5na^{t-2}-\frac{2a^{t-1}}t-\frac{4a^{t-2}}t\\
  &\ge2\left(1-\frac1t\right)e
       +\frac{a^{t-2}(n-3)}t(2a)
       +5na^{t-2}-\frac{2a^{t-1}}t-\frac{4a^{t-2}}t\\
  &=2e+\frac{a^{t-2}}t\bigl((5t-6)n+2\bigr)>2e+6.
\end{split}                                                   \tag{AST3.18}
\]
The last excess is at least $3na^{t-2}\ge630$ for
$n\ge15$, $a\ge14$, $t\ge3$.  It contradicts (AST3.11),
proving the theorem for every tensor count in its stated range.
\end{proof}

\paragraph{Checks, sources and remaining scope.}
The all-size arguments are universal proofs.  A separate exact
supplement checks weighted cyclic spectra, the full rotation error,
actual letter-action indices, pair selection and the symbolic sum;
it does not enumerate Hurwitz components or infer an infinite
theorem from finite data.  Known normalizer, scalar-field and
actual-cover inputs are the complete reviewed D-AST-01--03 and
HD-AST-04--09.  The elementary cyclotomic input in HD-AST-05 is
Milne, \emph{Class Field Theory}, v4.03, Chapter V, Example 3.2,
printed p.~157.\cite{HD:MilneCFT2020}  The weighted-cycle proof
is credited to D-ATT-08, with its exact short-orbit defect retained.

This result restricts actual finite local systems and possible
period-rank-nine suppliers.  It supplies no algebraic cycle and
makes no new positivity or motivic assertion.  It does not decide
the Eisenstein case with $t\ge3$, nontransitive tensor configurations,
other ordinary alternating labels or faithful spin modules.  In
particular the original field is never silently replaced by the
larger inducing character field.
